% Source: Wald GR, physical PDF pp. 170--171
%         (printed pp. 157--158).
% Coverage: Chapter 6 problems 1--6 and end material.
% Figures/tables/footnotes: no figures, tables, or footnotes.
% Status: source-reviewed and compile-clean.
% Uncertainties: none.

\subsection*{Problems}

\begin{enumerate}[label=\arabic*., leftmargin=*, itemsep=1.1em]
\item Let \(M\) be a three-dimensional manifold possessing a spherically
  symmetric Riemannian metric with \(\nabla_a r\ne0\), where \(r\) is defined by
  Eq.~\eqref{eq:6.1.3}.
  \begin{enumerate}[label=\alph*), leftmargin=2.2em, itemsep=.45em]
  \item Show that a new isotropic radial coordinate \(\bar r\) can be introduced
    so that the metric takes the form
    \[
      ds^2=H(\bar r)[d\bar r^2+\bar r^2d\Omega^2].
    \]
    (This shows that every spherically symmetric three-dimensional space is
    conformally flat.)
  \item Show that in isotropic coordinates the Schwarzschild metric is
    \[
      ds^2=-\frac{(1-M/2\bar r)^2}{(1+M/2\bar r)^2}\,dt^2+
      \left(1+\frac{M}{2\bar r}\right)^4
      [d\bar r^2+\bar r^2d\Omega^2].
    \]
  \end{enumerate}

\item Calculate the Ricci tensor, \(R_{ab}\), for a static, spherically symmetric
  spacetime, Eq.~\eqref{eq:6.1.5}, using the coordinate component method of
  Section~3.4a. Compare the amount of labor involved with that of the tetrad
  approach given in the text.

\item Consider the source-free (\(j^a=0\)) Maxwell's equations
  \eqref{eq:4.3.12} and \eqref{eq:4.3.13} in a static, spherically symmetric
  spacetime, Eq.~\eqref{eq:6.1.5}.
  \begin{enumerate}[label=\alph*), leftmargin=2.2em, itemsep=.45em]
  \item Argue that the general form of a Maxwell tensor which shares the static
    and spherical symmetries of the spacetime is
    \[
      F_{ab}=2A(r)(e_0)_{[a}(e_1)_{b]}
      +2B(r)(e_2)_{[a}(e_3)_{b]},
    \]
    where the \((e_\mu)_a\) are defined by Eq.~\eqref{eq:6.1.6a}.
  \item Show that if \(B(r)=0\), the general solution of Maxwell's equations with
    the form of part (a) is \(A(r)=-q/r^2\), where \(q\) may be interpreted as the
    total charge. [The solution obtained with \(B(r)\ne0\) is a duality rotation
    of this solution, representing the field of a magnetic monopole.]
  \item Write down and solve Einstein's equation, \(G_{ab}=8\pi T_{ab}\), with
    electromagnetic stress-energy tensor corresponding to the solution of part
    (b). Show that the general solution is the Reissner--Nordstr\"om metric,
    \[
      ds^2=-\left(1-\frac{2M}{r}+\frac{q^2}{r^2}\right)dt^2+
      \left(1-\frac{2M}{r}+\frac{q^2}{r^2}\right)^{-1}dr^2+r^2d\Omega^2.
    \]
  \end{enumerate}

\item Let \((M,g_{ab})\) be a stationary spacetime with timelike Killing field
  \(\xi^a\). Let \(V^2=-\xi^a\xi_a\).
  \begin{enumerate}[label=\alph*), leftmargin=2.2em, itemsep=.45em]
  \item Show that the acceleration
    \(a^b=u^a\nabla_a u^b\) of a stationary observer is given by
    \(a^b=\nabla^b\ln V\).
  \item Suppose in addition that \((M,g_{ab})\) is asymptotically flat, i.e., that
    there exist coordinates \(t,x,y,z\) [with \(\xi^a=(\partial/\partial t)^a\)]
    such that the components of \(g_{ab}\) approach
    \(\operatorname{diag}(-1,1,1,1)\) as \(r\to\infty\), where
    \(r=(x^2+y^2+z^2)^{1/2}\). (See Chapter~11 for further discussion of
    asymptotic flatness.) As in the case of the Schwarzschild metric, the energy
    as measured at infinity of a particle of mass \(m\) and 4-velocity \(u^a\)
    is \(E=-m\xi^au_a\). Suppose a particle of mass \(m\) is held stationary by a
    massless string, with the other end of the string held by a stationary observer
    at large \(r\). Let \(F\) denote the force exerted by the string on the
    particle. According to part (a) we have
    \(F=mV^{-1}[\nabla_aV\nabla^aV]^{1/2}\). Use conservation-of-energy
    arguments to show that the force exerted by the observer at infinity on the
    other end of the string is \(F_\infty=VF\). Thus, the magnitude of the force
    exerted at infinity differs from the force exerted locally by the redshift
    factor.
  \end{enumerate}

\item Derive the formula, Eq.~\eqref{eq:6.3.45}, for the general relativistic time
  delay.

\item Show that any particle (not necessarily in geodesic motion) in region II
  (\(r<2M\)) of the extended Schwarzschild spacetime, Figure~\ref{fig:6.9}, must
  decrease its radial coordinate at a rate given by
  \[
    \left|dr/d\tau\right|\geq[2M/r-1]^{1/2}.
  \]
  Hence, show that the maximum lifetime of any observer in region II is
  \(\tau=\pi M\,[\sim10^{-5}(M/M_\odot)\,\mathrm{s}]\), i.e., any observer in
  region II will be pulled into the singularity at \(r=0\) within this proper
  time. Show that this maximum time is approached by freely falling, i.e.,
  geodesic, motion from \(r=2M\) with \(E\to0\).
\end{enumerate}
